\documentclass[10pt,a4paper]{article}
\usepackage{fullpage}
\usepackage[utf8]{inputenc}
\usepackage[T1]{fontenc}
\usepackage{amsmath}
\usepackage{amssymb}
\usepackage{graphicx}
\usepackage[english]{babel}
\usepackage{url}
\usepackage{hyperref}
\usepackage{listings}
\usepackage{textcomp}
\usepackage{accsupp}
\usepackage{attachfile}
\usepackage{verbatim}
\usepackage[misc]{ifsym}
\usepackage{placeins}


% color def
\usepackage{color}
\definecolor{darkred}{rgb}{0.6,0.0,0.0}
\definecolor{darkgreen}{rgb}{0,0.50,0}
\definecolor{lightblue}{rgb}{0.0,0.42,0.91}
\definecolor{orange}{rgb}{0.99,0.48,0.13}
\definecolor{grass}{rgb}{0.18,0.80,0.18}
\definecolor{pink}{rgb}{0.97,0.15,0.45}

% listings
\usepackage{listings}

% General Setting of listings
\lstset{
	language=Python,
	tabsize=4,
	aboveskip=1em,
	breaklines=true,
	abovecaptionskip=-6pt,
	captionpos=b,
	escapeinside={\%*}{*)},
	frame=single,
	numbers=left,
	numbersep=8pt,
	numberstyle=\tiny,
	morekeywords=[1]{,as,assert,nonlocal,with,yield,self,True,False,None,} % Python builtin
	morekeywords=[2]{,__init__,__add__,__mul__,__div__,__sub__,__call__,__getitem__,__setitem__,__eq__,__ne__,__nonzero__,__rmul__,__radd__,__repr__,__str__,__get__,__truediv__,__pow__,__name__,__future__,__all__,}, % magic methods
	morekeywords=[3]{,object,type,isinstance,copy,deepcopy,zip,enumerate,reversed,list,set,len,dict,tuple,range,xrange,append,execfile,real,imag,reduce,str,repr,}, % common functions
	morekeywords=[4]{,Exception,NameError,IndexError,SyntaxError,TypeError,ValueError,OverflowError,ZeroDivisionError,}, % errors
	morekeywords=[5]{,ode,fsolve,sqrt,exp,sin,cos,arctan,arctan2,arccos,pi, array,norm,solve,dot,arange,isscalar,max,sum,flatten,shape,reshape,find,any,all,abs,plot,linspace,legend,quad,polyval,polyfit,hstack,concatenate,vstack,column_stack,empty,zeros,ones,rand,vander,grid,pcolor,eig,eigs,eigvals,svd,qr,tan,det,logspace,roll,min,mean,cumsum,cumprod,diff,vectorize,lstsq,cla,eye,xlabel,ylabel,squeeze,}, % numpy / math
	deletekeywords=[2]{next,set},
	deletekeywords=[3]{next,set},
	basicstyle=\ttfamily,
	backgroundcolor=\color{white},
	commentstyle=\color{darkgreen}\slshape,
	keywordstyle=\color{blue}\bfseries\itshape,
	keywordstyle=[2]\color{blue}\bfseries,
	keywordstyle=[3]\color{grass},
	keywordstyle=[4]\color{red},
	keywordstyle=[5]\color{orange},
	stringstyle=\color{darkred},
	emphstyle=\color{pink}\underbar,
	morestring=[d]{"},
	showstringspaces=false,
	upquote=true,
    literate={á}{{\'a}}1
        	{é}{{\'e}}1
        	{í}{{\'{\i}}}1
        	{ó}{{\'o}}1
        	{ú}{{\'u}}1
        	{Á}{{\'A}}1
        	{É}{{\'E}}1
        	{Í}{{\'I}}1
        	{Ó}{{\'O}}1
        	{Ú}{{\'U}}1
        	{ü}{{\"u}}1
        	{Ü}{{\"U}}1
        	{ñ}{{\~n}}1
        	{Ñ}{{\~N}}1
        	{¿}{{?`}}1
        	{¡}{{!`}}1
}

\usepackage{tikz}
\usetikzlibrary{shapes.geometric, arrows}
\tikzstyle{startstop} = [rectangle, rounded corners, minimum width=3cm, minimum height=1cm,text centered, draw=black, fill=red!30]
\tikzstyle{io} = [trapezium, trapezium left angle=70, trapezium right angle=110, minimum width=3cm, minimum height=1cm, text centered, draw=black, fill=blue!30]
\tikzstyle{process} = [rectangle, minimum width=3cm, minimum height=1cm, text centered, draw=black, fill=orange!30]
\tikzstyle{decision} = [diamond, minimum width=3cm, minimum height=1cm, text centered, draw=black, fill=green!30]
\tikzstyle{arrow} = [thick,->,>=stealth]


\makeatletter
\lst@RequireAspects{writefile}

% Use \attachfile command to add listing content as paperclip.
\lstnewenvironment{attachedlisting}[1][]{%
	\lst@BeginAlsoWriteFile{#1.txt}%
}
{%
	\lst@EndWriteFile
	\marginpar{\attachfile[appearance=false,icon=Paperclip,mimetype=text/tex]{#1.txt}}
}

% Use accsupp package to add listing content as copyable text.
\lstnewenvironment{copyablelisting}{%
	\lst@BeginAlsoWriteFile{\jobname.txt}%
}
{%
	\lst@EndWriteFile
	\let\verbatim@processline\add@lstline
	\global\let\lstfile\empty
	\verbatiminput{\jobname.txt}%
	\marginpar{(\BeginAccSupp{method=escape,ActualText={\lstfile}}\PaperPortrait\EndAccSupp{})}
}
\def\add@lstline
{\xdef\lstfile{\unexpanded\expandafter{\lstfile}\the\verbatim@line\string^^J}}
\makeatother


\title{{\huge Verificación de Paréntesis Balanceados}}
\date{}

\begin{document}
\maketitle


\section{Objetivo}
Dada una cadena de expresión, \lstinline|expression|, escribe un programa para examinar si los pares y el orden de ``\{'', ``\}'', ``('', ``)'', ``['', ``]'', ``$\langle$'', ``$\rangle$'' son correctos en \lstinline|expression|. Por ejemplo:
\begin{description}
\item[Balanceado:] \lstinline|expresion = "[()]{}{()()}"|
\item[No balanceado:] \lstinline|expresion = "[(])"|
\end{description}

El programa debe funcionar con cualquier tipo de cadena que contenga cualquier carácter, no solo paréntesis o corchetes.

\section{Algoritmo}
\begin{enumerate}
\item Crear una Pila.
\item Recorrer la cadena de expresión
    \begin{itemize}
    \item Si el carácter actual es un paréntesis de apertura (``('' o ``\{'' o ``['' o ``$\langle$'') entonces empújalo a la pila
    \item Si el carácter actual es un paréntesis de cierre (``)'' o ``\}'' o ``]'' o ``$\rangle$''), entonces saca de la pila y verifica que el carácter sacado sea el paréntesis de apertura correspondiente. Si lo es, entonces está balanceado y continúa verificando, de lo contrario, no está balanceado. Ten en cuenta que la pila puede o no tener elementos restantes.
    \end{itemize}
\item Después de recorrer completamente, si queda algún paréntesis de apertura en la pila, entonces la expresión no está balanceada.
\end{enumerate}

\section{Funcionalidad a Implementar}

El programa mostrará un menú con dos opciones para verificar si la expresión de entrada está balanceada:

\begin{description}
\item[Verificar una cadena introducida por el usuario: ] El usuario introducirá directamente la expresión usando la consola.
\item[Verificar el contenido de un archivo: ] El usuario elegirá un archivo usando un selector de archivos, lee la Sección \ref{sec:easy}. El programa usará el contenido de todo el archivo como la expresión.
\end{description}

\subsection{Ejemplo de Ejecución del Programa}

El Listing~\ref{lst:sample} contiene un ejemplo de la ejecución del programa.

\begin{lstlisting}[abovecaptionskip=0pt, caption={Ejemplo de ejecución del programa.}, label=lst:sample]
#***********************************
#1.-> Verificar balanceo de paréntesis (introducir texto en pantalla)
#2.-> Verificar balanceo de paréntesis (seleccionar archivo)
#3.-> Salir
#***********************************
#Elige una opción: 1
#Por favor introduce el texto a verificar: (8+(3+4)+ 5 + [3*4]
#No balanceado
#
#***********************************
#1.-> Verificar balanceo de paréntesis (introducir texto en pantalla)
#2.-> Verificar balanceo de paréntesis (seleccionar archivo)
#3.-> Salir
#***********************************
#Elige una opción: 1
#Por favor introduce el texto a verificar: (())[]
#Balanceado
#
#***********************************
#1.-> Verificar balanceo de paréntesis (introducir texto en pantalla)
#2.-> Verificar balanceo de paréntesis (seleccionar archivo)
#3.-> Salir
#***********************************
#Elige una opción: 1
#Por favor introduce el texto a verificar: [()]{}{()()}
#Balanceado
#
#***********************************
#1.-> Verificar balanceo de paréntesis (introducir texto en pantalla)
#2.-> Verificar balanceo de paréntesis (seleccionar archivo)
#3.-> Salir
#***********************************
#Elige una opción: 1
#Por favor introduce el texto a verificar: [()}
#No balanceado
#
#***********************************
#1.-> Verificar balanceo de paréntesis (introducir texto en pantalla)
#2.-> Verificar balanceo de paréntesis (seleccionar archivo)
#3.-> Salir
#***********************************
#Elige una opción: 3
#Saliendo....
\end{lstlisting}

\section{Otras Observaciones}

\begin{itemize}
\item Usa funciones para estructurar tu programa y minimizar el uso de código fuera de las funciones.
    \begin{itemize}
    \item Usa la función main como el punto de entrada del programa.
\begin{attachedlisting}[main]
if __name__ == "__main__":
    main()
\end{attachedlisting}
    \end{itemize}
\item Desafíos adicionales:
    \begin{itemize}
    \item Intenta usar solo las ventanas creadas por EasyGui para interactuar con el usuario, es decir, crea un programa que solo utilice una interfaz gráfica.
    \item Puedes agregar opciones para filtrar tipos de archivos según las extensiones de archivo. Consulta la documentación de \texttt{fileopenbox} y los parámetros que toma.
    \item Si la expresión no está balanceada, ¿puedes decirle al usuario en qué punto se rompe la condición? ¿Y dónde está el paréntesis correspondiente? Intenta usar \texttt{colorama} para pintar la parte correcta en verde y la incorrecta en rojo.
    \end{itemize}
\end{itemize}

\section{EasyGui}
\label{sec:easy}

Para implementar la segunda opción, tendremos que importar una librería adicional de Python llamada \href{https://easygui.readthedocs.io/en/latest/api.html}{EasyGui}\footnote{\href{https://easygui.readthedocs.io/en/latest/api.html}{https://easygui.readthedocs.io/en/latest/api.html}}. EasyGui proporciona una interfaz fácil de usar para la interacción simple de GUI (Interfaz Gráfica de Usuario) con un usuario. No requiere que el programador sepa nada sobre elementos de GUI más complejos como tkinter, frames, widgets, callbacks o lambdas. Todas las interacciones de GUI se invocan mediante llamadas a funciones simples que devuelven resultados.

Para usar la librería, primero tenemos que instalarla en nuestras computadoras. Para hacer esto, podemos usar la consola que proporciona Spyder. Usaremos la utilidad \lstinline|pip| para instalar las librerías necesarias, en sus computadoras, dependiendo de su instalación de Python, es posible que tengan que usar \lstinline|conda|, el comando debería ser el mismo, simplemente cambien una palabra por la otra.

\begin{lstlisting}
pip install easygui
\end{lstlisting}

Si la instalación es exitosa, Spyder nos dirá que es posible que necesitemos reiniciar el kernel. Esto se hace haciendo clic derecho en el nombre de la pestaña de la consola en la parte superior derecha de la consola y eligiendo la opción de reiniciar el kernel.

La librería EasyGui proporciona varias funciones para abrir diferentes tipos de ventanas, para seleccionar archivos, para solicitar la entrada del usuario, para mostrar un mensaje, etc. Aquí hay algunos ejemplos de cómo usar algunas de esas funciones, si deseas más información, puedes acceder a la documentación de la librería siguiendo el enlace. Recuerda importar siempre la librería.

\subsection{Selector de Archivos}
Una de las opciones, y la que tendrás que elegir para completar este proyecto de laboratorio, es una ventana para elegir un archivo. La función toma varios parámetros, todos opcionales, y devuelve el nombre del archivo que el usuario ha elegido.

\begin{attachedlisting}[open]
# Abrir un selector de archivos con un nombre para la ventana
nombre = easygui.fileopenbox(title="Verificador de paréntesis")
# Intenta ver cómo se comporta el nombre del archivo
# ¿Qué pasa si el usuario no elige un archivo?
print(nombre)

# Abrir un selector de archivos que solo muestra archivos .txt
nombre = easygui.fileopenbox(default=["*.txt"])
# ¿Podemos limitar a otros tipos de archivos?
# ¿Debería estar balanceado un archivo fuente de Python?
\end{attachedlisting}

\subsection{Mostrar un Mensaje}
Podemos mostrar al usuario una ventana con un mensaje, cuando el usuario presione OK, la ventana se cerrará.

\begin{attachedlisting}[message]
# Mostrar un mensaje simple
easygui.msgbox(msg="Mensaje")

# Podemos agregar un título
easygui.msgbox(msg="Mensaje", title="Error")

# Incluso podemos modificar lo que se muestra en el botón de ok
easygui.msgbox(msg="Mensaje", title="Error", ok_button="Soy un botón")
\end{attachedlisting}

\subsection{Opción Sí/No}
Una interacción más compleja sería hacerle una pregunta Sí/No al usuario.

\begin{attachedlisting}[yesno]
# Una pregunta simple al usuario
eleccion = easygui.ynbox(msg="¿Intentar de nuevo?")
# ¿Cómo sabemos qué ha elegido el usuario?
# ¿Cuáles son los posibles valores para esta variable?
print(eleccion)

# También podemos modificar lo que muestran los botones
eleccion = easygui.ynbox(msg="¿Intentar de nuevo?", choices=["Sí", "No"])
# ¿Esto cambia lo que se devuelve?
\end{attachedlisting}

\subsection{Mostrar Múltiples Opciones}
Incluso podemos pedirle al usuario que seleccione una de un conjunto de opciones que le damos. La siguiente ventana mostrará al usuario las opciones que queremos y una vez que haya elegido una de ellas, sabremos el índice de la elegida.

\begin{attachedlisting}[choices]
# Primero definimos una lista de opciones
opciones = ["primera", "segunda", "tercera"]
# Ahora creamos una ventana con esas opciones
eleccion = easygui.choicebox(choices=opciones)
# ¿Cuáles son los posibles valores para esta variable?
print(eleccion)

# Podemos usar esto para crear el menú de nuestro programa
eleccion = easygui.choicebox(msg="Bienvenido al menú principal\nSelecciona una opción", title="Verificador de paréntesis", choices=opciones)
\end{attachedlisting}

Si las posibles opciones no son tantas, podemos mostrarlas usando un conjunto de botones.

\begin{attachedlisting}[buttons]
# Establecemos las opciones como antes
opciones = ["primera", "segunda", "tercera"]
eleccion = easygui.buttonbox(choices=opciones)
# ¿Cómo sabemos qué botón ha presionado el usuario?
# No tenemos un índice como antes
print(eleccion)
\end{attachedlisting}

\subsection{Solicitar Cualquier Entrada}
También podemos pedirle al usuario que escriba algo, de manera similar a la función \lstinline|input()| de Python.

\begin{attachedlisting}[input]
# Simplemente le pediremos al usuario que ingrese algo
respuesta = easygui.enterbox(msg="Escribe algo")
# ¿Qué pasa si el usuario presiona cancelar o cierra la ventana?
print(respuesta)

# A veces queremos que el usuario tenga más espacio
# Podemos usar una función diferente
respuesta = easygui.textbox(msg="Escribe mucho")
# ¿Qué pasa si el usuario presiona enter?
# ¿Qué ha pasado en el enterbox si el usuario presionó enter?
\end{attachedlisting}

\subsection{Y Aún Más}
Intenta ver qué sucede al ejecutar los fragmentos de código mostrados aquí. La mayoría de las funciones aceptarán \lstinline|msg| y \lstinline|title| como parámetros. Hay más alternativas incluidas en esta librería, puedes usar el siguiente fragmento de código para ver cómo se verían. Si deseas más información, puedes consultar la documentación de la librería.

\begin{attachedlisting}[demo]
easygui.egdemo()
\end{attachedlisting}

\newpage

\end{document}